% Appendix
% Note: \appendix command is in main.tex

\section{Training Details}
\label{app:training_details}

Table~\ref{tab:hyperparams} lists the complete hyperparameter settings for all experimental conditions.

\begin{table}[h]
\centering
\small
\caption{Hyperparameters for all experimental conditions.}
\label{tab:hyperparams}
\resizebox{\columnwidth}{!}{%
\begin{tabular}{@{}lcc@{}}
\toprule
Parameter & LoRA & Full-FT \\
\midrule
Base model & \multicolumn{2}{c}{Qwen3-0.6B} \\
Reward model & \multicolumn{2}{c}{ArmoRM-Llama3-8B-v0.1} \\
Prompts (train / eval) & \multicolumn{2}{c}{2000 / 200} \\
Generations & \multicolumn{2}{c}{3} \\
Temperature & \multicolumn{2}{c}{0.7} \\
Max new tokens & \multicolumn{2}{c}{1024} \\
Epochs per generation & 2 & 2 (C4, C8: 3) \\
Per-device batch size & 4 & 1 \\
Gradient accumulation & 4 & 8 \\
Optimizer & \multicolumn{2}{c}{AdamW} \\
Weight decay & \multicolumn{2}{c}{0.01} \\
Warmup ratio & \multicolumn{2}{c}{0.1} \\
Precision & \multicolumn{2}{c}{bf16} \\
Learning rate & $2 \times 10^{-4}$ & $5 \times 10^{-6}$ / $2 \times 10^{-5}$ \\
LoRA rank & 16 & -- \\
LoRA alpha & 32 & -- \\
LoRA target & \parbox{3.2cm}{\raggedright q/k/v/o\_proj,\\gate/up/down\_proj} & -- \\
Gradient checkpointing & No & Yes \\
\bottomrule
\end{tabular}}
\end{table}

\section{LoRA as Capacity Bottleneck}
\label{app:lora_bottleneck}

The similarity of diversity outcomes between $K{=}4$ and $K{=}8$ LoRA conditions (C1: $-$0.27\%, C2: $+$0.04\%) despite doubled selection pressure suggests that parameter capacity, not curation intensity, is the binding constraint under LoRA. If curation were the primary driver of diversity change, doubling the candidate pool from 4 to 8 (reducing the retention rate from 25\% to 12.5\%) should produce measurably greater diversity suppression. Instead, both conditions converge to nearly identical Gen~2 Vendi Scores (107.25 for both C1 and C2). This convergence is consistent with a capacity-limited regime in which the adapter cannot express the differences between curated and non-curated training signals.

The rank-16 LoRA adapter modifies only 1.67\% of model parameters. All learned changes, whether driven by reward-curated data or random data, must be expressed within this low-dimensional subspace. Training loss provides direct evidence of this constraint: C1 (curated, $K{=}4$) reaches a final loss of 0.166, C2 (curated, $K{=}8$) reaches 0.168, and C5 (random) reaches 0.152. The 0.014 gap between curated and random final loss is small relative to the 0.240 reduction from Gen~0 to Gen~2 within each condition (Table~\ref{tab:training_loss}). The adapter absorbs most of the training signal's structure regardless of whether that signal comes from curated or random selection; the residual difference in loss translates to a negligible difference in diversity outcomes.

This interpretation reframes the 3.40~pp curation gap reported in Section~\ref{sec:direction_reversal}. The gap does not reflect curation actively suppressing diversity. Rather, random data's uncorrelated gradients happen to push the LoRA subspace in directions that slightly increase output diversity, while curated data's correlated gradients (consistent reward-preferred patterns) do not. Under full fine-tuning, where the capacity constraint is removed, this asymmetry disappears and eventually reverses (Section~\ref{sec:direction_reversal}). The practical implication is that LoRA-based studies of reward curation effects should be interpreted cautiously. Observed diversity changes under parameter-efficient fine-tuning may reflect adapter capacity limits rather than properties of the curation strategy itself.

\section{Per-Generation Metrics}
\label{app:per_gen_metrics}

Table~\ref{tab:per_gen_metrics} reports four diversity and fluency metrics for every condition at each self-training generation.
Vendi Score measures semantic diversity via matrix-rank entropy over embedding similarities; higher values indicate a more diverse output distribution.
Self-BLEU averages pairwise BLEU between generated responses; lower values indicate less repetition.
Perplexity (PPL) is computed with a held-out language model; stable values indicate no fluency degradation.
Nearest-neighbor distance (NND) is the mean cosine distance to the closest other response in embedding space; higher values reflect greater spread.

\begin{table*}[t]
\centering
\small
\caption{Per-generation metrics for all 8 experimental conditions across 3 self-training generations. Curated conditions (C1--C4, C8) show stable Vendi Score ($\pm$1\%), low Self-BLEU, and consistent NND, indicating no semantic collapse. Random baselines (C5--C7) exhibit comparable or larger metric drift despite the absence of reward selection. C7 metrics were recorded from a separate analysis pass.}
\label{tab:per_gen_metrics}
\begin{tabular}{@{}llcccc@{}}
\toprule
Condition & Gen & Vendi Score & Self-BLEU ($\downarrow$) & PPL & NND \\
\midrule
\multicolumn{6}{l}{\textit{Curated conditions}} \\
\midrule
C1: LoRA, $K{=}4$
  & 0 & 107.54 & 0.0059 & 10.60 & 0.7088 \\
  & 1 & 106.73 & 0.0091 &  9.73 & 0.7065 \\
  & 2 & 107.25 & 0.0062 &  9.72 & 0.7054 \\
\midrule
C2: LoRA, $K{=}8$
  & 0 & 107.21 & 0.0057 &  9.40 & 0.7097 \\
  & 1 & 107.26 & 0.0024 &  9.18 & 0.7086 \\
  & 2 & 107.25 & 0.0064 & 10.20 & 0.7095 \\
\midrule
C3: Full-FT, 5e-6
  & 0 & 107.10 & 0.0038 &  9.10 & 0.7073 \\
  & 1 & 107.61 & 0.0060 &  8.66 & 0.7091 \\
  & 2 & 108.07 & 0.0044 &  9.21 & 0.7108 \\
\midrule
C4: Full-FT, 2e-5
  & 0 & 107.81 & 0.0046 &  9.57 & 0.7079 \\
  & 1 & 108.03 & 0.0039 & 10.34 & 0.7107 \\
  & 2 & 108.54 & 0.0040 &  8.82 & 0.7100 \\
\midrule
C8: Full-FT, $K{=}16$
  & 0 & 107.32 & 0.0057 &  8.50 & 0.7057 \\
  & 1 & 106.73 & 0.0050 &  8.02 & 0.7065 \\
  & 2 & 107.11 & 0.0058 &  7.71 & 0.7092 \\
\midrule
\multicolumn{6}{l}{\textit{Random baselines}} \\
\midrule
C5: LoRA, $K{=}4$
  & 0 & 105.42 & 0.0044 &  9.19 & 0.7005 \\
  & 1 & 108.21 & 0.0067 &  8.76 & 0.7121 \\
  & 2 & 108.72 & 0.0067 & 11.18 & 0.7136 \\
\midrule
C6: Full-FT, 5e-6
  & 0 & 107.17 & 0.0045 &  8.77 & 0.7061 \\
  & 1 & 107.87 & 0.0047 &  8.70 & 0.7094 \\
  & 2 & 108.64 & 0.0065 &  8.82 & 0.7114 \\
\midrule
C7: Full-FT, 2e-5
  & 0 & 108.08 & 0.0046 & 10.97 & 0.7064 \\
  & 1 & 107.85 & 0.0049 & 10.83 & 0.7095 \\
  & 2 & 108.13 & 0.0068 & 11.15 & 0.7094 \\
\bottomrule
\end{tabular}
\end{table*}

\section{ArmoRM Quality Scores}
\label{app:quality_scores}

Table~\ref{tab:quality_scores} reports mean ArmoRM quality scores per generation for all conditions, computed on 200 held-out evaluation prompts.

\begin{table}[h]
\centering
\small
\caption{Mean ArmoRM quality score per generation. All conditions remain in the 0.110--0.116 range with mild decreases across generations ($-$0.3\% to $-$1.3\%), confirming quality stability under iterative self-training.}
\label{tab:quality_scores}
\begin{tabular}{@{}llcccc@{}}
\toprule
& Config & Gen~0 & Gen~1 & Gen~2 & $\Delta$ \\
\midrule
\multicolumn{6}{@{}l}{\textit{Curated}} \\
C1 & LoRA, $K{=}4$ & 0.1133 & 0.1132 & 0.1142 & $+$0.08\% \\
C2 & LoRA, $K{=}8$ & 0.1146 & 0.1136 & 0.1143 & $-$0.03\% \\
C3 & Full-FT, 5e-6 & 0.1145 & 0.1141 & 0.1131 & $-$1.2\% \\
C4 & Full-FT, 2e-5 & 0.1130 & 0.1154 & 0.1124 & $-$0.5\% \\
C8 & Full-FT, $K{=}16$ & 0.1162 & 0.1162 & 0.1149 & $-$1.1\% \\
\midrule
\multicolumn{6}{@{}l}{\textit{Random baseline}} \\
C5 & LoRA, $K{=}4$ & 0.1107 & 0.1099 & 0.1096 & $-$1.0\% \\
C6 & Full-FT, 5e-6 & 0.1126 & 0.1111 & 0.1117 & $-$0.8\% \\
C7 & Full-FT, 2e-5 & 0.1112 & 0.1100 & 0.1097 & $-$1.3\% \\
\bottomrule
\end{tabular}
\end{table}

\section{Lexical Diversity: Self-BLEU and Distinct-$n$}
\label{app:lexical_diversity}

Table~\ref{tab:lexical_diversity} reports Self-BLEU-4 (computed over 500 randomly sampled responses per generation, seed 42) and Distinct-$n$ ($n \in \{1,2,3\}$, computed over all 2{,}000 training responses). Distinct-$n$ is defined as the ratio of unique $n$-grams to total $n$-grams across the response corpus.

\begin{table*}[t]
\centering
\small
\caption{Self-BLEU-4 and Distinct-$n$ across generations. Self-BLEU-4 remains low ($<$0.016) for all conditions, indicating minimal textual repetition. Distinct-$n$ shows uniform decline across all conditions (curated and random alike), consistent with format homogenization under iterative self-training rather than reward-specific narrowing.}
\label{tab:lexical_diversity}
\begin{tabular}{@{}llccccc@{}}
\toprule
Condition & Gen & Self-BLEU-4 ($\downarrow$) & Distinct-1 & Distinct-2 & Distinct-3 \\
\midrule
\multicolumn{6}{l}{\textit{Curated conditions}} \\
\midrule
C1: LoRA, $K{=}4$
  & 0 & 0.0116 & 0.167 & 0.582 & 0.835 \\
  & 1 & 0.0160 & 0.163 & 0.560 & 0.808 \\
  & 2 & 0.0137 & 0.157 & 0.546 & 0.793 \\
\midrule
C2: LoRA, $K{=}8$
  & 0 & 0.0138 & 0.167 & 0.579 & 0.831 \\
  & 1 & 0.0144 & 0.164 & 0.562 & 0.810 \\
  & 2 & 0.0149 & 0.162 & 0.555 & 0.802 \\
\midrule
C3: Full-FT, 5e-6
  & 0 & 0.0106 & 0.166 & 0.579 & 0.832 \\
  & 1 & 0.0144 & 0.165 & 0.569 & 0.819 \\
  & 2 & 0.0158 & 0.162 & 0.556 & 0.803 \\
\midrule
C4: Full-FT, 2e-5
  & 0 & 0.0131 & 0.168 & 0.580 & 0.833 \\
  & 1 & 0.0117 & 0.162 & 0.562 & 0.811 \\
  & 2 & 0.0101 & 0.160 & 0.557 & 0.807 \\
\midrule
C8: Full-FT, $K{=}16$
  & 0 & 0.0120 & 0.170 & 0.581 & 0.832 \\
  & 1 & 0.0156 & 0.165 & 0.570 & 0.819 \\
  & 2 & 0.0160 & 0.163 & 0.562 & 0.811 \\
\midrule
\multicolumn{6}{l}{\textit{Random baselines}} \\
\midrule
C5: LoRA, $K{=}4$
  & 0 & 0.0107 & 0.168 & 0.580 & 0.831 \\
  & 1 & 0.0136 & 0.158 & 0.545 & 0.789 \\
  & 2 & 0.0132 & 0.156 & 0.534 & 0.776 \\
\midrule
C6: Full-FT, 5e-6
  & 0 & 0.0133 & 0.169 & 0.582 & 0.833 \\
  & 1 & 0.0133 & 0.165 & 0.563 & 0.811 \\
  & 2 & 0.0121 & 0.162 & 0.550 & 0.795 \\
\midrule
C7: Full-FT, 2e-5
  & 0 & 0.0133 & 0.167 & 0.580 & 0.832 \\
  & 1 & 0.0138 & 0.162 & 0.556 & 0.800 \\
  & 2 & 0.0141 & 0.157 & 0.541 & 0.783 \\
\bottomrule
\end{tabular}
\end{table*}

